\subsection{Sposoby oceny jakości automatycznego code review}
\label{subsec:llm_acr_eval}

Właściwe mierzenie skuteczności systemów do automatycznego przeglądu kodu jest problematyczną kwestią.
Największym problemem jest to, że klasyczne miary używane do przetwarzania języka naturalnego rzadko kiedy
oddają to, czy komentarz code review jest faktycznie przydatny technicznie \cite{Lin2025CodeReviewQA, Zhang2024SurveyLLMSE}.
W publikacjach widać trzy główne podejścia do ocen: metryki ilościowe, badania jakościowe z użyciem programistów oraz
gotowe benchmarki badawcze.

\paragraph{Metryki ilościowe i ich ograniczenia}
W zadaniach typu \textit{code-to-comment} czy przy automatycznych poprawkach (\textit{code refinement}), badacze zazwyczaj sięgają po 
metryki leksykalne - mowa tu głównie o BLEU-4, ROUGE-L oraz METEOR \cite{Lin2025CodeReviewQA, Meng2025RARe, Yu2024Carllm}. 
Mimo ich powszechnego użycia, wiarygodność tych miar jest coraz częściej podważana w nowszych pracach:

\begin{itemize}
    \item \textbf{Krytyka Exact Match (EM):} Ta metryka bywa oceniana jako zbyt restrykcyjna, 
zwłaszcza w zadaniach, gdzie liczy się sens, a nie identyczne brzmienie. 
W pracy Tufano \cite{Tufano2024CodeReviewAutomation} ręczna analiza pokazała, że nawet 21,83\% predykcji odrzuconych przez 
EM w zadaniu generowania komentarzy było jednak semantycznie poprawnych. To dość mocno sugeruje, że samo EM może 
"karać" poprawne odpowiedzi tylko dlatego, że są sformułowane inaczej. 
W pracy Pornprasit \cite{Pornprasit2024FineTuningPromptingCR} pojawia się podobny wniosek, gdzie zwraca się uwagę, że 
EM może zaniżać realne możliwości modeli w porównaniu z oceną ludzką (która jest bardziej elastyczna).
    
    \item \textbf{BERTScore:} Jest to podejście oparte o aspekt kontekstowy, więc zamiast patrzeć 
na identyczność tokenów, próbuje uchwycić podobieństwo znaczenia. Dla komentarzy do code review ma to sens, bo pozwala lepiej 
oddać intencję recenzenta niż proste dopasowanie słów, co podkreślają autorzy pracy Haider \cite{Haider2024PromptingFineTuningRCG}.
    
    \item \textbf{Desiredness Score (DS):} Opiera się na różnicy \textit{perplexity} modelu.
Zamysł jest taki, że mierzy się, na ile dany komentarz zmniejsza niepewność modelu przy generowaniu poprawnej wersji kodu.
Dzięki temu da się automatycznie wyłapywać tzw. \textit{Desired Review Comments} (DRC), czyli takie uwagi, które faktycznie prowadzą do 
napraw, a nie tylko "brzmią sensownie" \cite{Yu2025Desiview}.
    
    \item \textbf{Wskaźniki poprawności technicznej:} Coraz częściej podkreśla się potrzebę mierzenia nie tylko 
jakości językowej komentarzy, ale też efektu w kodzie. W szczególności istotne są \textit{Correction Ratio} (czy poprawki realnie działają)
oraz \textit{Regression Ratio}, które opisuje ryzyko, że sugestia modelu popsuje wcześniej poprawny
kod \cite{Cihan2025EvaluatingLLMsForCodeReview}. To jest o tyle ważne, że nawet "ładna" uwaga w code review nie ma wartości,
jeśli w praktyce prowadzi do regresji.
\end{itemize}

\paragraph{Oceny eksperckie i badania użyteczności}
Podczas typowego wdrożenia przemysłowego automatyczne metrki to z reguły za mało, więc dalej kluczową rolę w procesach
code review odgrywa przegląd przez doświadczonego programistę \cite{Ramesh2025EricssonExperience}. 
W nowych badaniach dodatkowo zaczęto odchodzić od ogólnych ocen na rzecz wielowymiarowych skal.
W takim podejściu zwraca się uwagę na trafność (\textit{relevance}), czyli to, czy dany komentarz w 
ogóle odnosi się do właściwego fragmentu kodu i czy poprawnie wskazuje źródło problemu \cite{Yu2024Carllm, Ren2025HydraReviewer}.
Kolejnym jest użyteczność (\textit{actionability}) - czy sugestia jest na tyle konkretna, żeby programista wiedział, jak ma wprowadzić poprawkę.
Jest to o tyle istotne, że według badań Ren \cite{Ren2025HydraReviewer} aż 61,0\% uwag generowanych przez tradycyjne metody 
jest ocenianych jako zbyt ogólnikowe (\textit{vagueness}), co czyni je bezużytecznymi.
Ponad to jest kwestia zrozumiałości (\textit{comprehensibility}) - czy wyjaśnienie błędu jest jasne dla człowieka i
czy model, generując odpowiedź, zachował logiczny ciąg rozumowania, co często wiąże się z zastosowaniem
techniki \textit{Chain-of-Thought} \cite{Yu2024Carllm, Haider2024PromptingFineTuningRCG}.

Istnieje także podejście \textit{LLM-as-a-judge}. Jest to wykorzystanie zaawansowanych i większych modeli do oceny
wyników pracy mniejszych systemów. Wyniki takich eksperymentów wyglądają obiecująco — w pracy Jaoua \cite{Jaoua2025LLMStaticAnalyzersCR} wykazano, 
że taka automatyczna ocena wykazuje dość wysoką zgodność z opiniami ludzi, osiągając współczynnik \textit{Cohen's 
kappa} na poziomie 0,72.

\paragraph{Benchmarki i dekompozycja zadań}
Przegląd kodu co raz rzadziej jest traktowany jako jedno, monolityczne zadanie. Częściej 
rozbija się je na etapy odpowiadające kolejnym krokom rozumowania. Przykładem jest benchmark \textit{CodeReviewQA}, w 
którym sprawdza się działanie modelu w trzech krokach \cite{Lin2025CodeReviewQA}:
\begin{enumerate}
    \item \textbf{CTR (\textit{Change Type Recognition}):} rozpoznanie, jaki typ zmiany jest potrzebny.
    \item \textbf{CL (\textit{Change Localisation}):} wskazanie konkretnych linii kodu, których dotyczy problem. W 
praktyce to właśnie lokalizacja bywa najczęstszym "wąskim gardłem" dla modeli LLM \cite{Lin2025CodeReviewQA}.
    \item \textbf{SI (\textit{Solution Identification}):} wybór właściwej poprawki spośród dostępnych alternatyw.
\end{enumerate}

Trzeba dodatkowo brać pod uwagę problemy metodologiczne. Dane historyczne 
potrafią być zaszumione: w popularnych zbiorach szum szacuje się na około 25\% rekordów \cite{Tufano2024CodeReviewAutomation}.
Istnieje ryzyko \textit{data leakage} (wycieku danych), które w zadaniach naprawy kodu 
potrafi zawyżać wyniki o ponad 30 punktów procentowych \cite{Germano2025VulnSLR}.
